% GENERATED FILE. Edit canonical lessons, then run npm run generate:books.
% canonical-source-hash: fnv1a64:5b8e8ce6807df723
% canonical-lessons: GE-C375-klausur, GE-C375-durchfallen, GE-C375-muendlich, GE-C375-schriftlich, GE-C375-zeugnis

\chapter{Die Klausur}
\label{ch:ge-die-klausur}
\hlchaptermodality{375}

\begin{chapteropening}
\textbf{By the end of this chapter:} \emph{I can talk about written and oral exams, failing one, and a school report.}

Complete the last lesson of chapter 375: I can talk about written and oral exams, failing one, and a school report.
\end{chapteropening}

\section[die Klausur]{die Klausur — the written exam}

\label{lesson:GE-C375-klausur}



\begin{quote}
Before the new one: say the German for flexible, then the German for the planning.
\end{quote}

\subsection*{You'll want to know: die Klausur}
\textbf{die Klausur} — \textquotedblleft{}the written exam\textquotedblright{}.

\begin{quote}
Morgen schreiben wir eine Klausur in Mathe.
\end{quote}

\begin{quote}
\emph{Tomorrow we are sitting a written exam in maths.}
\end{quote}

Plural: \emph{die Klausuren}. At school in Austria the written test is \emph{die Schularbeit}.

\begin{cousinweb}
\emph{Klausur} is a learned borrowing from Latin \emph{clausūra} \textquotedblleft{}a lock, a bar\textquotedblright{} (Lewis \& Short), from \emph{claudere} \textquotedblleft{}to close\textquotedblright{}: you sit the exam shut away from help.
\end{cousinweb}

\begin{grammarlens}[title={What you've built}]
One more word for telling, explaining and describing at B1.
\end{grammarlens}

\subsection*{Your turn}
\begin{itemize}
\raggedright
  \item \emph{Say it:} \emph{die Klausur}
  \item \emph{Say it:} \emph{die Klausur}, once more
  \item \emph{Say it:} a sentence of your own with \emph{die Klausur}
  \item \emph{Recall:} say the German for flexible, then the German for the planning, then say \emph{die Klausur} again
\end{itemize}

\begin{tcolorbox}[breakable,colback=teal!4,colframe=teal!35!black,title={Before you move on}]
What does die Klausur mean? (\textquotedblleft{}The written exam\textquotedblright{}.) Say it once more.
\end{tcolorbox}
\section[durchfallen]{durchfallen — to fail an exam}

\label{lesson:GE-C375-durchfallen}



\begin{quote}
Before the new one: say the German for the planning, then the German for the written exam.
\end{quote}

\subsection*{You'll want to know: durchfallen}
\textbf{durchfallen} — \textquotedblleft{}to fail an exam\textquotedblright{}.

\begin{quote}
Mein Bruder ist bei der Fahrprüfung durchgefallen.
\end{quote}

\begin{quote}
\emph{My brother failed his driving test.}
\end{quote}

Separable, Perfekt with \emph{sein}: \emph{er fällt durch}, \emph{er ist durchgefallen}. \emph{Durch eine Prüfung fallen} — to fail an exam.

\begin{cousinweb}
\emph{durchfallen} is built from \emph{durch} \textquotedblleft{}through\textquotedblright{} and \emph{fallen} \textquotedblleft{}to fall\textquotedblright{}, inherited from Old High German \emph{fallan} and Proto-Germanic \emph{fallaną} (EtymDB): you fall through instead of getting past.
\end{cousinweb}

\begin{grammarlens}[title={What you've built}]
One more word for telling, explaining and describing at B1.
\end{grammarlens}

\subsection*{Your turn}
\begin{itemize}
\raggedright
  \item \emph{Say it:} \emph{durchfallen}
  \item \emph{Say it:} \emph{durchfallen}, once more
  \item \emph{Say it:} a sentence of your own with \emph{durchfallen}
  \item \emph{Recall:} say the German for the planning, then the German for the written exam, then say \emph{durchfallen} again
\end{itemize}

\begin{tcolorbox}[breakable,colback=teal!4,colframe=teal!35!black,title={Before you move on}]
What does durchfallen mean? (\textquotedblleft{}To fail an exam\textquotedblright{}.) Say it once more.
\end{tcolorbox}
\section[mündlich]{mündlich — oral, spoken}

\label{lesson:GE-C375-muendlich}



\begin{quote}
Before the new one: say the German for the written exam, then the German for to fail an exam.
\end{quote}

\subsection*{You'll want to know: mündlich}
\textbf{mündlich} — \textquotedblleft{}oral, spoken\textquotedblright{}.

\begin{quote}
Die mündliche Prüfung dauert fünfzehn Minuten.
\end{quote}

\begin{quote}
\emph{The oral exam lasts fifteen minutes.}
\end{quote}

\emph{Mündlich} comes from \emph{der Mund} — the mouth.

\begin{cousinweb}
\emph{mündlich} is built with \emph{-lich} on \emph{Mund} \textquotedblleft{}mouth\textquotedblright{}, inherited from Proto-Germanic \emph{munþaz} (EtymDB): by mouth.
\end{cousinweb}

\begin{grammarlens}[title={What you've built}]
One more word for telling, explaining and describing at B1.
\end{grammarlens}

\subsection*{Your turn}
\begin{itemize}
\raggedright
  \item \emph{Say it:} \emph{mündlich}
  \item \emph{Say it:} \emph{mündlich}, once more
  \item \emph{Say it:} a sentence of your own with \emph{mündlich}
  \item \emph{Recall:} say the German for the written exam, then the German for to fail an exam, then say \emph{mündlich} again
\end{itemize}

\begin{tcolorbox}[breakable,colback=teal!4,colframe=teal!35!black,title={Before you move on}]
What does mündlich mean? (\textquotedblleft{}Oral, spoken\textquotedblright{}.) Say it once more.
\end{tcolorbox}
\section[schriftlich]{schriftlich — written, in writing}

\label{lesson:GE-C375-schriftlich}



\begin{quote}
Before the new one: say the German for to fail an exam, then the German for oral.
\end{quote}

\subsection*{You'll want to know: schriftlich}
\textbf{schriftlich} — \textquotedblleft{}written, in writing\textquotedblright{}.

\begin{quote}
Die Kündigung muss schriftlich sein.
\end{quote}

\begin{quote}
\emph{The notice must be in writing.}
\end{quote}

The opposite of \emph{mündlich}. \emph{Schriftlich und mündlich} — written and spoken.

\begin{cousinweb}
\emph{schriftlich} is built with \emph{-lich} on \emph{Schrift} \textquotedblleft{}writing\textquotedblright{}, from Old High German \emph{scrift}, which goes back to Latin \emph{scrīptum} \textquotedblleft{}something written\textquotedblright{} (EtymDB).
\end{cousinweb}

\begin{grammarlens}[title={What you've built}]
One more word for telling, explaining and describing at B1.
\end{grammarlens}

\subsection*{Your turn}
\begin{itemize}
\raggedright
  \item \emph{Say it:} \emph{schriftlich}
  \item \emph{Say it:} \emph{schriftlich}, once more
  \item \emph{Say it:} a sentence of your own with \emph{schriftlich}
  \item \emph{Recall:} say the German for to fail an exam, then the German for oral, then say \emph{schriftlich} again
\end{itemize}

\begin{tcolorbox}[breakable,colback=teal!4,colframe=teal!35!black,title={Before you move on}]
What does schriftlich mean? (\textquotedblleft{}Written, in writing\textquotedblright{}.) Say it once more.
\end{tcolorbox}
\section[das Zeugnis]{das Zeugnis — the school report; the certificate, the reference}

\label{lesson:GE-C375-zeugnis}



\begin{quote}
Before the new one: say the German for oral, then the German for written.
\end{quote}

\subsection*{You'll want to know: das Zeugnis}
\textbf{das Zeugnis} — \textquotedblleft{}the school report; the certificate, the reference\textquotedblright{}.

\begin{quote}
Am Ende des Schuljahres bekommen die Kinder ihr Zeugnis.
\end{quote}

\begin{quote}
\emph{At the end of the school year the children get their reports.}
\end{quote}

Plural: \emph{die Zeugnisse}. From an employer: \emph{das Arbeitszeugnis} — a written reference.

\begin{cousinweb}
\emph{Zeugnis} is formed with \emph{-nis} on \emph{zeugen} \textquotedblleft{}to testify\textquotedblright{}, from Middle High German \emph{ziugen} (EtymDB): a paper that bears witness to what you did.
\end{cousinweb}

\begin{grammarlens}[title={What you've built}]
One more word for telling, explaining and describing at B1.
\end{grammarlens}

\subsection*{Your turn}
\begin{itemize}
\raggedright
  \item \emph{Say it:} \emph{das Zeugnis}
  \item \emph{Say it:} \emph{das Zeugnis}, once more
  \item \emph{Say it:} a sentence of your own with \emph{das Zeugnis}
  \item \emph{Recall:} say the German for oral, then the German for written, then say \emph{das Zeugnis} again
\end{itemize}

\begin{tcolorbox}[breakable,colback=teal!4,colframe=teal!35!black,title={Before you move on}]
What does das Zeugnis mean? (\textquotedblleft{}The school report; the certificate, the reference\textquotedblright{}.) Say it once more.
\end{tcolorbox}
